%Tex root = ProbabilitàEStatistica.tex

\chapter{Indici di posizione e dispersione}
Al giorno d'oggi non è raro dover trattare grandi quantità di dati.\\[0.2ex]
Per avere una sensazione di un così vasto campione di dati, è utile saperli sintettizzare in qualche misura

\section{Indici di posizione}
Gli \textcolor{viola}{\textbf{indici di posizione}} \hl{aiutano a capire quali sono i valori che assume una certa distribuzione} con particolare attenzione al centro di questa.

\subsection{Media}
Consideriamo di avere la seguente distribuzione: N rappresenta la numerisità della popolazione, n il numero delle classi.

\begin{center}
    \begin{tabular}{|c|c|c|}
        \hline
        \textbf{X} & \textbf{Frequenze} & \textbf{Frequenze relative}\\
        \hline
        $\mathrm{x_1}$ & $\mathrm{f_1}$ & $\mathrm{p_1}$\\
        $\mathrm{x_2}$ & $\mathrm{f_2}$ & $\mathrm{p_2}$\\
        $\vdots$ & $\vdots$ & $\vdots$\\
        $\mathrm{x_n}$ & $\mathrm{f_n}$ & $\mathrm{p_n}$\\
        \hline
        \textcolor{red}{\textbf{Totale}} & \textcolor{red}{\textbf{N}} & \textcolor{red}{\textbf{1}}\\
        \hline
    \end{tabular}
\end{center}

\Definizione{media}{
Si chiama \textbf{media} di X è la quantità:
\begin{equation*}
    \boldsymbol{\mathrm{\bar{X} =\frac{1}{N}\sum_{i=1}^{n}x_i \cdot f_i = \sum_{i=1}^{n}xi \cdot p_i}}
\end{equation*}
}

\esempio{calcolo della media}{
Consideriamo la distribuzione di alcuni studenti.

\smallskip
\begin{center}
    \begin{tabular}{|c|c|c|}
    \hline
    \textbf{Peso (Kg)} & \textbf{Frequenze} & \textbf{Frequenze relative}\\
    \hline
    65 & 2 & 0.2\\
    68 & 1 & 0.1\\
    69 & 1 & 0.1\\
    70 & 2 & 0.2\\
    72 & 1 & 0.1\\
    74 & 3 & 0.3\\
    \hline
    \textcolor{red}{\textbf{Totale}} & \textcolor{red}{\textbf{10}} & \textcolor{red}{\textbf{1.0}}\\
    \hline
    \end{tabular}
\end{center}

\medskip
Si ha:
\begin{equation*}
    \boldsymbol{\mathrm{\bar{P}}} = \frac{1}{10}(65 \cdot 2 + 68 \cdot 1 + 69 \cdot 1 + 70 \cdot 2 + 72 \cdot 1 + 74 \cdot 3) = \boldsymbol{70.1}
\end{equation*}
}

\subsection{Mediana, quatrili e percentili}
\Definizione{mediana}{
La \textbf{mediana} è il valore che provoca la ripartione della popolazione in esame in due parti ugulamente numerose: per il 50$\%$ della popolazione il dato è minore della mediana, per il restante 50$\%$ il dato è maggiore della mediana.
}

Se abbiamo i dati non raggruppati, possiamo pensarli sotto forma di fila ordinata.\\[0.2ex]
La mediana, allora, è:
\begin{itemize}[label=$\triangleright$]
\item \textbf{valore centrale}: numeri dispari
\smallskip
\item \textbf{media dei valori centrali}: numeri pari
\end{itemize}

\esempio{calcolo della mediana}{
Consideriamo il seguente campione X di numeri:
\smallskip
\begin{center}
    67, 72, 78, 78, 84, 85, 87, 91 
\end{center}

\smallskip
Si tratta di un campione di numerosità 8 (pari):
\begin{equation*}
    \mathrm{\boldsymbol{Med(X)}} = \frac{1}{2}(78 + 84) = 81
\end{equation*}

\bigskip
Consideriamo il seguente campione X di numeri:
\begin{center}
    65, 67, 72, 78, 78, 84, 85, 87, 91 
\end{center}

\smallskip
Si tratta di un campione di numerosità 9 (dipari):
\begin{equation*}
    \mathrm{\boldsymbol{Med(X)}} = 78
\end{equation*}
}

La \hl{media e la mediana sono statitsiche utili per descrivere i valori centrali dei dati}.\\[0.2ex]
La media fa uso di tutti i dati e, in particolare, è influenzata in maniera sensibile dai valori eccezionalmente alti o bassi.\\[0.2ex]
La mediana, invece, dipende direttamente solo da uno o due valori in centro alla distribuzione e non risente dei dati estremi.

\bigskip
\Definizione{quartili}{
I \textbf{quartili} sono quei valori che ripartiscono la popolazione (ordinata), in quattro parti ugulamente numerose (ciasuna pari al 25$\%$ del totale):
\begin{itemize}[label=$\diamond$]
    \item \textbf{primo quartile ($q_1$)}: 25$\%$ della popolazione a sinistra e il 75$\%$ a destra
    \smallskip
    \item \textbf{secondo quartile ($q_2$)}: 50$\%$ della popolazione a sinistra e il 50$\%$ a destra (mediana)
    \smallskip
    \item \textbf{terzo quartile ($q_3$)}: 75$\%$ della popolazione a sinistra e il 25$\%$ a destra
\end{itemize}
}

Se il campione si trova sotto forma di fila ordinata ($\mathrm{x_1, x_2, \dots, x_N}$), allora per il calcolo si procede in modo analogo a quanto fatto per la mediana.\\[0.2ex]
Se si vuole calcolare $\mathrm{q_k}$ si moltiplica $\mathrm{p = 0.k}$ per la numerosità del campione N:
\begin{itemize}
\item se pN è intero:
\begin{equation*}
    \mathrm{q_k = \frac{1}{2}(x_{pN} + x_{pN + 1})}
\end{equation*}
\smallskip
\item se pN non è un numero intero:
\begin{equation*}
    \mathrm{q_k = x_{\lfloor pN \rfloor + 1}}
\end{equation*}
\end{itemize}

\esempio{calcolo dell'n-simo quartile}{
Consideriamo il seguente campione X di numeri:
\smallskip
\begin{center}
    67, 72, 78, 78, 84, 85, 87, 91 
\end{center}

\smallskip
    Si tratta di un campione di numerosità 8 (pari). 
    
    \textcolor{red}{Quanto valgono i quartili?}
    \begin{itemize}[label=$\triangleright$]
    \item \textbf{primo quartile}: 0.25 $\cdot$ 8 = 2
    \begin{equation*}
    \mathrm{q_1 = \frac{1}{2}(x_2 + x_3)} = \frac{1}{2}(72+78) = \boldsymbol{75}
    \end{equation*}
    \item \textbf{secondo quatrile}:
    \begin{equation*}
    \mathrm{q_2 = Med(X) = \frac{1}{2}(78 + 84)} = \boldsymbol{81}
    \end{equation*}
    \item \textbf{terzo quatrile}: 0.75 $\cdot$ 8 = 6
    \begin{equation*}
    \mathrm{q_3 = \frac{1}{2}(x_6 + x_7)} = \frac{1}{2}(85 + 87) = \boldsymbol{86}
    \end{equation*}
    \end{itemize}
}

\subsection{Moda}

\Definizione{moda}{
La \textbf{moda} di un insieme di dati, se esiste, è l'\hl{unico valore che ha la freuqunza massima}.

Se non vi è solo valore con frequenza massima, ciascuno di essi è detto \textbf{valore modale}.
}

\textcolor{red}{Perchè sccegliere la moda, invece, di media o mediana?}\\[0.2ex]
La moda identifica il picco della distribuzione (scelta prevalente), mentre la media può rappresentare un valore che non è mai stato osservato.


\section{Indici di dispersione}
Gli \textcolor{viola}{\textbf{indici di disperisione }} indicano quanto i dati sono concetrati o viceversa dispersi attorno a tali valori tipici.

\subsection{Varianza e scarto quadratico medio}
È interessante misurare, quindi, il grado di disperisione dei dati rispetto, ad esempio, alla media.\\[0.2ex]
Si osserva che la somma di tutte le derivazioni dalla media è sempre zero:
\begin{equation*}
\mathrm{\sum_{i=1}^{n}(x_i - \bar{X}) \cdot f_i = 0}
\end{equation*} 

Perciò, per misurare in modo significativo la disperisione dei dati rispetto alla media, si può considerare, ad esempio, la somma dei moduli delle derivazioni, oppure la somma dei quadrati delle derivazioni.

\medskip
Consideriamo di avere la seguente distribuzione: N rappresenta la numerisità della popolazione.

\begin{center}
\begin{tabular}{|c|c|c|}
    \hline
    \textbf{X} & \textbf{Frequenze} & \textbf{Frequenze relative}\\
    \hline
    $\mathrm{x_1}$ & $\mathrm{f_1}$ & $\mathrm{p_1}$\\
    $\mathrm{x_2}$ & $\mathrm{f_2}$ & $\mathrm{p_2}$\\
    $\vdots$ & $\vdots$ & $\vdots$\\
    $\mathrm{x_n}$ & $\mathrm{f_n}$ & $\mathrm{p_n}$\\
    \hline
    \textcolor{red}{\textbf{Totale}} & \textcolor{red}{\textbf{N}} & \textcolor{red}{\textbf{1}}\\
    \hline
\end{tabular}
\end{center}

\Definizione{varianza}{
Si chiama \textbf{varianza} di X la media degli scarti al quadrato:
\begin{equation*}
    \boldsymbol{\mathrm{s_X^2 = \frac{1}{N}\sum_{i=1}^{n}(x_i - \bar{X})^2 \cdot f_i = \sum_{i=1}^{n}(x_i - \bar{X})^2 \cdot p_i}}
\end{equation*}
}


\Definizione{scarto quadratico medio}{
Si chiama \textbf{scarto quadratico medio} o \textbf{deviazione standard} di X la radice della varianza:
\begin{equation*}
    \boldsymbol{\mathrm{s_X = \sqrt{\frac{1}{N}\sum_{i=1}^{n}(x_i - \bar{X})^2 \cdot f_i} = \sqrt{\sum_{i=1}^{n}(x_i - \bar{X})^2 \cdot p_i}}}
\end{equation*}
}

La varianza e la deviazione standard sono grandezze non negative che si annullano solo quanto gli $\mathrm{x_i}$, sono tutti uguali tra loro e, quindi, uguali alla loro media.

\medskip
Il calcolo della varianza può essere semplificato se si nota che, prese comunque due costanti $\mathrm{a}$ e $\mathrm{b}$, se si considera il nuovo insieme di dati $\mathrm{y_i := ax_i + b}$, dove $\mathrm{i = 1, 2, \dots, n}$, allora per quanto già detto $\mathrm{\bar{Y} = a \bar{X} + b}$ e, quindi:
\begin{equation*}
\mathrm{\sum_{i=1}^{n}(y_i - \bar{Y})^2 = a^2 \sum_{i=1}^{n}(x_i - \bar{X})^2}
\end{equation*}


Perciò, se $\mathrm{s_y^2}$ e $\mathrm{s_x^2}$ sono le rispettive varianza campionarie, si ha che:
\begin{equation*}
\mathrm{s_y^2 = a^2 s_x^2}
\end{equation*}

\esempio{calcolo della varianza e deviazione standard}{
Consideriamo la seguente distribuzione di frequenze:

\smallskip
\begin{center}
    \begin{tabular}{|c|c|}
    \hline
    \textbf{Distrubuzione} & \textbf{Frequenze}\\
    \hline
    110 < D $\le$ 130 & 20\\
    130 < D $\le$ 150 & 40\\
    150 < D $\le$ 170 & 60\\
    170 < D $\le$ 210 & 80\\
    \hline
    \textcolor{red}{\textbf{Totoale}} & \textcolor{red}{\textbf{200}}\\
    \hline
    \end{tabular}
\end{center}

\bigskip
La media è uguale a:
\begin{equation*}
    \mathrm{\bar{D}} = \frac{1}{200}(120 \cdot 20 + 140 \cdot 40 + 160 \cdot 60 + 190 \cdot 80) = \boldsymbol{164}
\end{equation*}

\smallskip
\begin{center}
    \begin{tabular}{|c|c|c|c|}
    \hline
    \textbf{Distrubuzione} & $\boldsymbol{\mathrm{(D - \bar{D})^2}}$ & $\boldsymbol{\mathrm{D^2}}$ & \textbf{Frequenze}\\
    \hline
    110 < D $\le$ 130 & $(120 - 164)^2$ & $120^2$ & 20\\
    130 < D $\le$ 150 & $(140 - 164)^2$ & $140^2$ & 40\\
    150 < D $\le$ 170 & $(160 - 164)^2$ & $160^2$ & 60\\
    170 < D $\le$ 210 & $(190 - 164)^2$ & $190^2$ & 80\\
    \hline
    \textcolor{red}{\textbf{Totoale}} & & & \textcolor{red}{\textbf{200}}\\
    \hline
    \end{tabular}
\end{center}

\bigskip
La varianza è uguale a:
\begin{equation*}
    \mathrm{s_D^2} = \frac{1}{200}[(120 - 164)^2 \cdot 20 + (140 - 164)^2 \cdot 40 + (160 - 164)^2 \cdot 60 + (190 - 164)^2 \cdot 80] = \boldsymbol{584}
\end{equation*}

Lo scarto quadratico medio è uguale a:
\begin{equation*}
    \mathrm{ s_D = \sqrt{s_D^2}} = \boldsymbol{24.17}
\end{equation*}
}

\section{Campioni normali}
Osservando gli istogrammi dei campioni numerici forniti da esperimenti reali, si può notare che v sia una forma caratteristica che compare moloto spesso e accomuna un gran numero di campioni di dati, provenineti dai contesti più disparati.\\[0.2ex]
Questi grafici hanno un solo massimo, in corrispondenza della mediana, e decrescono da entrambi i lati simmestricamente, secondo una curva a campana.

%foto campioni normali

Campioni di dati che rispettano questi requeisiti si dice \textbf{normale}.

\bigskip
In realtà, pur mantendendo un aspetto simile a quello descritto, non capita mai che un istogramma reale rispetti perfettamente la simmetria e la monotonia.

%foto campioni swecked

Si può parlare, allora di campione apprositivamente normale.\\
Se un campione di dati presenta un istogramma che è sensibilmente asimmetrico rispetto alla mediana, si parla di \textbf{campione swecked} (sbilanciato) a sinistra o a destra, a seconda del lato in cui ha la più lunga.